\section*{Chapter \ref{ch:b1:earthquake-source} --- The Earthquake Source}

\begin{solution}{exo:b1:earthquake-source:1}
(a) $\log_{10}(\num{4e17}) = 17.60$, so $M_w = \tfrac23(17.60 - 9.1) = 5.67$.
(b) $\log_{10}(\num{5.59e22}) = 22.75$ and $M_w = \tfrac23(22.75 - 9.1) =
9.10$. (c) $M_0 = 10^{1.5\times7.0 + 9.1} = 10^{19.6} = \qty{3.98e19}{N.m}$.
\end{solution}

\begin{solution}{exo:b1:earthquake-source:2}
(a) $T = D/\dot s = \qty{3}{m}/\qty{0.025}{m.yr^{-1}} = \qty{120}{yr}$.
(b) With $c = 0.5$, $T = D/(c\dot s) = \qty{240}{yr}$: the creep takes up half
the motion, and the earthquakes, of the same slip, come half as often.
\end{solution}

\begin{solution}{exo:b1:earthquake-source:3}
By \cref{prop:b1:earthquake-source:moment-steps}, the ratio of the moments is
$10^{1.5\times(7.4 - 6.2)} = 10^{1.8} = 63.1$: it takes 63 earthquakes of
magnitude 6.2 to release the moment of one of magnitude 7.4.
\end{solution}

\begin{solution}{exo:b1:earthquake-source:4}
With $b = 1$, each unit of magnitude divides the count by 10: 10 earthquakes of
magnitude 5 or more a year, 1 of magnitude 6 or more, and 0.1 of magnitude 7 or
more, that is one every 10 years on average.
\end{solution}

\begin{solution}{exo:b1:earthquake-source:5}
(a) $A = \num{4.3e5} \times \qty{1.5e4}{m} = \qty{6.45e9}{m^2}$, and
$M_0 = \num{3e10}\times\num{6.45e9}\times4 = \qty{7.74e20}{N.m}$;
$\log_{10}M_0 = 20.89$ and $M_w = \tfrac23(20.89 - 9.1) = 7.86$. (b) A depth of
\qty{10}{km} divides the area and the moment by 1.5, which lowers the magnitude
by $\tfrac23\log_{10}1.5 = 0.12$, to 7.74. (c) The first, 7.86, rounds to the
7.9 of the catalogue; the second is a little low. The magnitude is not very
sensitive to the depth: a factor of 1.5 on the area is only 0.12 of magnitude.
\end{solution}

\begin{solution}{exo:b1:earthquake-source:6}
(a) $D = M_0/(\mu A) = \num{5.59e22}/(\num{3e10}\times\num{4e5}\times\num{1.5e5})
= \qty{31.1}{m}$. (b) $\qty{31.1}{m}/\qty{0.083}{m.yr^{-1}} = \qty{374}{yr}$.
(c) With the locked fraction 0.8, $31.1/(0.8\times0.083) = \qty{468}{yr}$; with
0.5, \qty{749}{yr}. The less the interface is locked, the longer the slip of
one great earthquake takes to build up; the average slip found here with a
rectangle of \qtyproduct{400 x 150}{km} and $\mu = \qty{30}{GPa}$ is itself
only an estimate, and the real slip was far from uniform.
\end{solution}

\begin{solution}{exo:b1:earthquake-source:7}
(a) The \omterm{def:b1:geological-maps:attitude}{strike} is $\phi = 329$, the \omterm{def:b1:geological-maps:attitude}{dip direction} $a = 059$, the \omterm{def:b1:geological-maps:attitude}{dip} $\delta =
\ang{8}$. Then $\mathbf s = (\sin329, \cos329, 0) = (-0.5150, 0.8572, 0)$,
$\mathbf u = (-\cos8\sin59, -\cos8\cos59, \sin8) = (-0.8488, -0.5100, 0.1392)$,
$\mathbf n = (\sin8\sin59, \sin8\cos59, \cos8) = (0.1193, 0.0717, 0.9903)$,
and with $r = \ang{110}$,
\[
  \mathbf d = \cos110\,\mathbf s + \sin110\,\mathbf u = (-0.6215, -0.7724, 0.1308) .
\]
(b) The \omterm{def:b1:earthquake-source:focal-mechanism}{auxiliary plane} has the normal $\mathbf d$ and the slip $\mathbf n$.
Taking the normal upward, it \omterm{def:b1:geological-maps:attitude}{dips} $\arccos 0.1308 = \ang{82.5}$ toward
$\arctan(-0.6215/-0.7724) = 218.8$, so its \omterm{def:b1:geological-maps:attitude}{strike} is $128.8$; the slip
$\mathbf n$ projected on its \omterm{def:b1:geological-maps:attitude}{strike} and \omterm{def:b1:geological-maps:attitude}{up-dip} vectors gives the \omterm{def:b1:faults-folds-fabrics:fault-elements}{rake}
\ang{87.2}. Rounded: $129/83/87$, the second plane of the published solution.
(c) $\mathbf P = (\mathbf n - \mathbf d)/\sqrt2$ \omterm{def:b1:field-lab-methods-1:trend-plunge}{trends} $221$ and \omterm{def:b1:field-lab-methods-1:trend-plunge}{plunges}
\ang{37}; $\mathbf T = (\mathbf n + \mathbf d)/\sqrt2$ \omterm{def:b1:field-lab-methods-1:trend-plunge}{trends} $036$ and \omterm{def:b1:field-lab-methods-1:trend-plunge}{plunges}
\ang{52} (published: $222/38$ and $36/52$). (d) The shallow plane, dipping
\ang{8} toward $059$, along the \omterm{def:b1:active-margins:slab}{subducting slab}: a \omterm{def:b1:faults-folds-fabrics:fault-types}{reverse fault} with a \omterm{def:b1:faults-folds-fabrics:fault-elements}{rake}
of \ang{110}, and, dipping less than \ang{45}, a thrust, the plate interface of
the subduction zone.
\end{solution}

\begin{solution}{exo:b1:earthquake-source:8}
(a) The plane \omterm{def:b1:geological-maps:attitude}{dips} toward $135$, so $\mathbf n = (\sin135, \cos135, 0) =
(0.707, -0.707, 0)$; a \omterm{def:b1:faults-folds-fabrics:fault-types}{dextral} slip has $r = \ang{180}$, $\mathbf d = -\mathbf s
= (-0.707, -0.707, 0)$. Toward the north, $\mathbf x = (0, 1, 0)$:
$(\mathbf n\cdot\mathbf x)(\mathbf d\cdot\mathbf x) = (-0.707)(-0.707) > 0$, a
push. East: $(0.707)(-0.707) < 0$, a pull. South: a push; west: a pull.
(b) $\mathbf P = (\mathbf n - \mathbf d)/\sqrt2 = (1, 0, 0)$: the \omterm{def:b1:earthquake-source:focal-mechanism}{P axis} is
horizontal and \omterm{def:b1:field-lab-methods-1:trend-plunge}{trends} $090$--$270$; $\mathbf T = (0, -1, 0)$ \omterm{def:b1:field-lab-methods-1:trend-plunge}{trends}
$000$--$180$. (c) The vertical plane striking $135$--$315$ has the normal
$\mathbf d$. With the slip $\mathbf n$, the block on its south-west side, into
which $\mathbf d$ points, moves toward $135$; seen from the north-east block,
looking south-west, that is to the left: a \omterm{def:b1:faults-folds-fabrics:fault-types}{sinistral} slip. By
\cref{prop:b1:earthquake-source:auxiliary}, this slip gives the same \omterm{def:b1:earthquake-source:focal-mechanism}{first
motions} in every direction.
\end{solution}

\begin{solution}{exo:b1:earthquake-source:9}
(a) $V = \qty{40}{mm.yr^{-1}}$ and $v = V/4$ at \qty{10}{km}, so by
\cref{prop:b1:earthquake-source:profile}(b) $D = \qty{10}{km}$; directly,
$D = x/\tan(\pi v/V) = 10/\tan(\pi/4) = \qty{10}{km}$. (b) $v(30) =
(40/\pi)\arctan 3 = \qty{15.9}{mm.yr^{-1}}$. (c) $V/(\pi D) =
\num{0.040}/(\pi\times\num{10000}) = \num{1.27e-6}$ per year.
\end{solution}

\begin{solution}{exo:b1:earthquake-source:10}
(a) The moments of the bins are still $C q^k$ with $q = 10^{1.5 - b}$, but now
$q < 1$: the moment of the bins decreases toward the large magnitudes, and the
bottom bin, $k = 0$, carries the share $C/S = (1 - q)/(1 - q^{K+1}) > 1 - q$, the
mirror image of \cref{prop:b1:earthquake-source:largest}. The smallest
earthquakes counted would then dominate, and the total would grow without limit
as smaller and smaller ones were included. (b) $1 - 10^{-0.7} = 0.80$: at least
\qty{80}{\percent}. (c) $\log_{10}N = 8.20 - 1.01\times9 = -0.89$, so $N = 0.13$
a year and about 6.4 in fifty years, three times the two observed. The line,
fitted between magnitudes 6 and 7.5, overestimates the largest earthquakes;
either the law bends down at the top, as the finite size of faults suggests, or
fifty years are too few to sample events that come once in decades. Both are
possible with two events.
\end{solution}

\begin{solution}{exo:b1:earthquake-source:11}
(a) In 150 years a post at $x$ moves by $150\,v(x) = 150\times(34/\pi)
\arctan(x/15)$ millimetres relative to the post on the fault: \qty{0.11}{m} at
\qty{1}{km}, \qty{1.28}{m} at \qty{15}{km}, \qty{2.08}{m} at \qty{50}{km},
toward the south for a \omterm{def:b1:faults-folds-fabrics:fault-types}{dextral} fault seen from the west. (b) The slip deficit is
$s = \qty{0.034}{m.yr^{-1}}\times\qty{150}{yr} = \qty{5.1}{m}$. During the
earthquake the posts move by $(s/\pi)\arctan(D/x)$: \qty{2.44}{m} at
\qty{1}{km}, \qty{1.28}{m} at \qty{15}{km}, \qty{0.47}{m} at \qty{50}{km}.
(c) The sums are $0.11 + 2.44$, $1.28 + 1.28$ and $2.08 + 0.47$, all
\qty{2.55}{m} $= s/2$: after the earthquake the fence is straight again on the
east side, displaced by $s/2$, and the west side, by symmetry, by $-s/2$.
\end{solution}

\begin{solution}{exo:b1:earthquake-source:12}
(a) $\mathbf P = (0, 1, 0)$ and $\mathbf T = (1, 0, 0)$. The pair
$\mathbf n = (\mathbf T + \mathbf P)/\sqrt2 = (0.707, 0.707, 0)$,
$\mathbf d = (\mathbf T - \mathbf P)/\sqrt2 = (0.707, -0.707, 0)$ satisfies
$\mathbf P = (\mathbf n - \mathbf d)/\sqrt2$ and $\mathbf T = (\mathbf n +
\mathbf d)/\sqrt2$. Both vectors are horizontal, so both \omterm{def:b1:earthquake-source:focal-mechanism}{nodal planes} are
vertical: the plane of normal $\mathbf n$ (toward $045$) \omterm{def:b1:geological-maps:attitude}{strikes} $135$--$315$,
the plane of normal $\mathbf d$ (toward $135$) \omterm{def:b1:geological-maps:attitude}{strikes} $045$--$225$. (b) On the
plane striking $315$, the block on the north-east side moves along $\mathbf d$,
toward $135$: seen from the other block, to the right, a \omterm{def:b1:faults-folds-fabrics:fault-types}{dextral} slip. On the
plane striking $045$, the block on the south-east side moves along $\mathbf n$,
toward $045$: a \omterm{def:b1:faults-folds-fabrics:fault-types}{sinistral} slip. (c) A new fault makes $\ang{45} - \phi/2 =
\ang{30}$ with $\sigma_1$: with $\sigma_1$ trending $000$, the faults would
\omterm{def:b1:geological-maps:attitude}{strike} $030$ (\omterm{def:b1:faults-folds-fabrics:fault-types}{sinistral}) and $330$ (\omterm{def:b1:faults-folds-fabrics:fault-types}{dextral}), \ang{15} from the \omterm{def:b1:earthquake-source:focal-mechanism}{nodal planes}.
The \omterm{def:b1:earthquake-source:focal-mechanism}{P axis} is fixed at \ang{45} from the \omterm{def:b1:earthquake-source:focal-mechanism}{nodal planes} by geometry, whatever the
friction of the rock, so it coincides with $\sigma_1$ only if the fault slipped
at \ang{45} from it. (d) The fault striking $040$ is \ang{5} from the \omterm{def:b1:earthquake-source:focal-mechanism}{nodal plane}
striking $045$: that plane is likely the fault, and the slip on it is
\omterm{def:b1:faults-folds-fabrics:fault-types}{sinistral}.
\end{solution}

\begin{solution}{pb:b1:earthquake-source:1}
\textbf{1.} The intervals are $1290 - 1050 = 240$, $1480 - 1290 = 190$ and
$1712 - 1480 = 232$ years; their mean is \qty{220.7}{yr}.

\textbf{2.} The mean offset is $(7.0 + 6.0 + 5.5 + 7.5)/4 = \qty{6.5}{m}$, and
$T = 6.5/0.030 = \qty{216.7}{yr}$, close to the mean interval.

\textbf{3.} The offsets of the last three earthquakes, $6.0 + 5.5 + 7.5 =
\qty{19.0}{m}$, accumulated between 1050 and 1712, in 662 years: $\qty{19.0}{m}
/\qty{662}{yr} = \qty{28.7}{mm.yr^{-1}}$, \qty{4}{\percent} less than the
geodetic \qty{30}{mm.yr^{-1}}. The two agree within the uncertainty of the
dates and offsets: the segment releases nearly all of its loading in its great
earthquakes.

\textbf{4.} $2025 - 1712 = 313$ years, and the deficit is
$\qty{0.030}{m.yr^{-1}}\times\qty{313}{yr} = \qty{9.39}{m}$.

\textbf{5.} The deficit exceeds every offset in the trench, the largest being
\qty{7.5}{m}, and the time since 1712 exceeds the longest interval, 240 years.
In that sense the segment is overdue; it does not mean that the earthquake is
imminent, only that the fault already holds more slip than it has ever
released at once in the record.

\textbf{6.} The intervals range from 190 to 240 years and the offsets from
\qtyrange{5.5}{7.5}{m}: the cycle is not periodic, and the strength of the
fault at the next earthquake is unknown. The trench gives a mean and a spread,
from which a probability can be estimated, not a date.

\textbf{7.} $A = \qty{250}{km}\times\qty{15}{km} = \qty{3750}{km^2} =
\qty{3.75e9}{m^2}$.

\textbf{8.} $M_0 = \num{3e10}\times\num{3.75e9}\times\num{9.39} =
\qty{1.06e21}{N.m}$.

\textbf{9.} $M_0 = \num{3e10}\times\num{3.75e9}\times\num{7.5} =
\qty{8.44e20}{N.m}$, so $M_w = \tfrac23(20.926 - 9.1) = 7.88$.

\textbf{10.} The ratio is $9.39/7.5 = 1.25$, which is $\tfrac23\log_{10}1.25 =
0.07$ of magnitude.

\textbf{11.} Half the area halves the moment and lowers the magnitude by
$\tfrac23\log_{10}2 = 0.20$: $M_w = 7.75$.

\textbf{12.} $\mu A V = \num{3e10}\times\num{3.75e9}\times\num{0.030} =
\qty{3.38e18}{N.m}$ per year.

\textbf{13.} $b = 1$. $N(5) = 10^{4.3 - 5} = 0.20$ a year and $N(6) = 0.020$ a
year.

\textbf{14.} $N(7) = 0.0020$ a year: one every 500 years; $N(8) = 0.00020$ a
year: one every \num{5000} years.

\textbf{15.} The law gives one earthquake of magnitude 7.95 or more every
$1/10^{4.3 - 7.95} = \num{4500}$ years, the trench one every 220 years, twenty
times more often. The law is fitted to the many small earthquakes of all the
faults of the region; extrapolated to the great earthquakes of one segment, it
underestimates them badly. The recurrence of great earthquakes on a fault is
taken from its own record and from its loading, not from the small ones.

\textbf{16.} $\num{5.5e17}/\num{3.38e18} = 0.16$: the moderate earthquakes
release \qty{16}{\percent} of the moment the segment accumulates.

\textbf{17.} By \cref{prop:b1:earthquake-source:moment-steps},
$10^{1.5\times(7.95 - 6)} = 10^{2.92} \approx 840$ times more.

\textbf{18.} With $b = 1$, the moment of the earthquakes is dominated by the
largest: whatever the moderate earthquakes do, they release a minor share, here
\qty{16}{\percent}, and the remaining \qty{84}{\percent}, the slip deficit,
waits for the great earthquake.

\textbf{19.} $v(-x) = -v(x)$ for each pair, so the profile is antisymmetric.
Far from the fault the velocities tend to $\pm V/2 = \pm\qty{15}{mm.yr^{-1}}$.

\textbf{20.} $D = x/\tan(\pi v/V) = 15/\tan(\pi\times7.5/30) = 15/\tan\ang{45} =
\qty{15}{km}$: at the \omterm{def:b1:earthquake-source:locking-depth}{locking depth} the velocity is $V/4$.

\textbf{21.} At \qty{5}{km}: $D = 5/\tan(\pi\times3.1/30) = 5/\tan\ang{18.6} =
\qty{14.9}{km}$; at \qty{40}{km}: $D = 40/\tan\ang{69.6} = \qty{14.9}{km}$. The
velocities are given to \qty{0.1}{mm.yr^{-1}}, and that rounding alone moves
$D$ by about \qty{0.1}{km}; with real data, the noise of each station would be
larger, and the \omterm{def:b1:earthquake-source:locking-depth}{locking depth} would be fitted to all the stations at once.

\textbf{22.} The \omterm{def:b1:earthquake-source:locking-depth}{locking depth}, \qty{15}{km}, equals the depth of the deepest
earthquakes. The locked part of the fault is its brittle part, the \omterm{def:b1:stress-strain-rheology:strength-envelope}{seismogenic
layer}, where friction holds it between earthquakes and fails in them; below, the
warmer rock creeps (\cref{ch:b1:stress-strain-rheology}) and the fault slides
steadily, which is the slip at the rate $V$ below $D$ of the model.

\textbf{23.} $V/(\pi D) = \num{0.030}/(\pi\times\num{15000}) = \num{6.4e-7}$
per year.

\textbf{24.} With $M_0 = \qty{1.056e21}{N.m}$, $\log_{10}M_0 = 21.024$ and
$M_w = \tfrac23(21.024 - 9.1) = 7.949$. To 0.01, the earthquake that would
release the whole deficit of the segment has the \omterm{def:b1:earthquake-source:mw}{moment magnitude}
\textbf{7.95}.
\end{solution}
